\chapter{Un même bord : flux intérieur et tension de jauge}
\label{chap:hoja-comun-ns-ym}
\index{Navier--Stokes}\index{Yang--Mills}
\index{opérateur de bord}\index{Stokes}

L'hydrodynamique et Yang--Mills ne sont pas réunis en comparant de l'extérieur deux
équations célèbres. Ils proviennent de deux représentations du même opérateur de
bord. L'orientation que le lecteur intérieur interprète comme circulation
devient, sous le lecteur réfléchi, une holonomie ; la mémoire symétrique
qui contrôle le flux devient un noyau positif de transfert.

\section{L'état commun de bord}

Soit \(K_m\) un complexe fini orienté, \(B_m\) son algèbre de données de
bord et \(d_m\) son cobord. L'état enrichi est
\[
 s_m=(b_m,U_m,\theta_m,c_m,\Omega_m),
\]
où \(U_m\) transporte, \(\theta_m^2=I\) inverse l'orientation,
\(c_m\) conserve la mémoire et \(\Omega_m\) est une \(1\)-cochaîne de connexion.
Les projecteurs de parité sont
\[
 E_\pm=\frac12(I\pm\theta_m).
\]
La partie \((U_m+U_m^\dagger)/2\) est symétrique ; ce n'est un projecteur que si l'on
démontre en outre son idempotence.

Sur le complexe de cochaînes, on construit deux lecteurs linéaires :
\[
 \rho_{{\rm int},m}:C^0(K_m)\longrightarrow
 C^1(K_m)\oplus C^0(K_m),
\qquad
 \rho_{{\rm int},m}(v)=(d_mv,d_m^\ast d_mv),
\]
\[
 \rho_{{\rm ref},m}:C^1(K_m)\longrightarrow
 C^0(K_m)\oplus C^1(K_m),
\qquad
 \rho_{{\rm ref},m}(a)=(d_m^\ast a,d_md_m^\ast a).
\]
Le premier transporte un potentiel de sommets vers son flux d'arêtes et sa
réponse intérieure ; le second transporte une cochaîne d'arêtes vers sa
réponse de bord et son opérateur réfléchi. Les codomaines sont distincts,
mais tous deux sont construits avec le même cobord. Circulation, vorticité,
holonomie et courbure apparaissent ensuite au moyen de l'identité de Stokes ; elles ne
sont pas présupposées à l'intérieur des symboles \(\rho\).

\section{L'opérateur minimal et son spectre}

Dans le cycle orienté à trois sommets, soit
\[
 d_0=
 \begin{pmatrix}
 -1&1&0\\
 0&-1&1\\
 1&0&-1
 \end{pmatrix}.
\]
Les laplaciens de Hodge de sommets et d'arêtes sont
\[
 \mathcal L_{\rm H}^{(0)}=d_0^\ast d_0,
 \qquad
 \mathcal L_{\rm H}^{(1)}=d_0d_0^\ast,
\]
et un calcul direct donne
\[
 \mathcal L_{\rm H}^{(0)}=\mathcal L_{\rm H}^{(1)}=
 \begin{pmatrix}
 2&-1&-1\\
 -1&2&-1\\
 -1&-1&2
 \end{pmatrix}
 =3I-J.
\]
Dans cet emplacement, les exposants \((0)\) et \((1)\) désignent les
laplaciens de Hodge sur les sommets et les arêtes. Ils ne désignent pas les opérateurs de
relèvement \(L_0,L_1\) de la chaîne \(w_6\to w_{30}\). Pour
\(\mathcal L_{\rm cyc}^{\rm TPK}:=2I-C-C^{-1}=3I-J\), la réalisation de
l'opérateur cyclique du TPK est typée au moyen de
\[
 \mathfrak R_{\rm hyd}(\mathcal L_{\rm cyc}^{\rm TPK})
 =\mathcal L_{\rm H}^{(0)},\qquad
 \mathfrak R_{\rm YM}(\mathcal L_{\rm cyc}^{\rm TPK})
 =\mathcal L_{\rm H}^{(1)}.
\]

\begin{theorem}[Opérateur commun de bord]
\label{thm:operador-comun-frontera-ns-ym}
Dans le cycle à trois cellules,
\[
 \operatorname{Spec}(\mathcal L_{\rm H}^{(0)})
 =\operatorname{Spec}(\mathcal L_{\rm H}^{(1)})
 =\{0,3,3\}.
\]
Après retrait du mode constant ---moyenne nulle dans le lecteur intérieur et
jauge dans le lecteur de connexion--- les deux secteurs possèdent un écart spectral exact \(3\).
Dans tout complexe fini, \(d^\ast d\) et \(dd^\ast\) ont le même spectre
non nul, avec multiplicités.
\end{theorem}

\begin{proof}
Le vecteur \((1,1,1)\) engendre le noyau de \(3I-J\). Sur son complément
orthogonal, \(J=0\), de sorte que l'opérateur est \(3I\). Pour l'énoncé
général, une décomposition en valeurs singulières \(d=U\Sigma V^\ast\) produit
\[
 d^\ast d=V\Sigma^\ast\Sigma V^\ast,
 \qquad
 dd^\ast=U\Sigma\Sigma^\ast U^\ast.
\]
Les carrés des valeurs singulières non nulles, y compris leurs
multiplicités, coïncident.
\end{proof}

\section{Lecture hydrodynamique}

Nous définissons le pas implicite
\[
 S_{\rm NS}=(I+\mathcal L_{\rm H}^{(0)})^{-1}.
\]
Dans le cycle minimal,
\[
 S_{\rm NS}=\frac14(I+J).
\]

\begin{theorem}[Contraction globale du flux fini]
\label{thm:contraccion-flujo-finito}
Pour tout \(v_0\) de moyenne nulle, l'itération
\(v_{k+1}=S_{\rm NS}v_k\) existe et est unique pour tout \(k\geq0\), et
\[
 v_k=4^{-k}v_0,
 \qquad
 \|v_k\|\leq4^{-k}\|v_0\|,
 \qquad
 \|v_k\|^2=16^{-k}\|v_0\|^2.
\]
\end{theorem}

\begin{proof}
À moyenne nulle, on a \(Jv=0\), donc
\(S_{\rm NS}v=v/4\). La formule s'obtient par récurrence.
\end{proof}

\Needspace{4\baselineskip}
Une non-linéarité orientée \(N_m(v,c)\) conserve l'énergie lorsque
\[
 \langle v,N_m(v,c)\rangle=0.
\]
L'advection redistribue alors l'énergie et l'opérateur positif contrôle
sa propagation. La mémoire \(c_m\) n'est pas absorbée dans une viscosité choisie :
elle fait partie de l'état qui doit être transporté entre les échelles.

\begingroup
\setlength{\parskip}{0pt}
\titlespacing*{\section}{0pt}{8pt plus 1pt minus .5pt}{2.5pt}
\setlength{\abovedisplayskip}{.25\baselineskip}
\setlength{\belowdisplayskip}{.25\baselineskip}
\setlength{\abovedisplayshortskip}{.12\baselineskip}
\setlength{\belowdisplayshortskip}{.18\baselineskip}
\section{Lecture Yang--Mills}

Une \(1\)-cochaîne \(a\) est interprétée comme un potentiel de connexion et possède
une énergie quadratique
\[
 Q_{\rm YM}(a)=\langle a,\mathcal L_{\rm H}^{(1)}a\rangle.
\]

\begin{theorem}[Coercivité et décroissance de jauge finies]
\label{thm:coercividad-gap-gauge-finito}
Dans le quotient par le mode de jauge constant,
\[
 Q_{\rm YM}(a)\geq3\|a\|^2.
\]
Le semi-groupe \(K_t=e^{-t\mathcal L_{\rm H}^{(1)}}\) satisfait
\[
 \|K_ta\|\leq e^{-3t}\|a\|
\]
dans le secteur physique du cycle minimal.
\end{theorem}

\begin{proof}
Sur le complément du noyau, le
\cref{thm:operador-comun-frontera-ns-ym} donne
\(\mathcal L_{\rm H}^{(1)}=3I\). Les deux formules
découlent du calcul fonctionnel.
\end{proof}

Pour une connexion non abélienne,
\[
 F_\Omega=d\Omega+\Omega\smile\Omega,
\qquad
 W_\gamma(\Omega)=\operatorname{Tr}\operatorname{Hol}_\gamma(\Omega).
\]
La positivité par réflexion requiert une factorisation sur la moitié
positive,
\[
 \langle\theta f,Kf\rangle=\|Cf\|^2\geq0;
\]
la positivité ponctuelle de la matrice \(K\) ne la remplace pas.

\section[Identité de Stokes : circulation--vorticité et holonomie--courbure]{Identité discrète de Stokes et \(2\) réalisations : circulation--vorticité et holonomie--courbure}

Pour une \(1\)-cochaîne \(\Omega\) et une \(2\)-chaîne \(K\),
\[
 \sum_{e\subset\partial K}\Omega(e)
 =
 \sum_{f\in K^{(2)}}(d\Omega)(f).
\]
Les arêtes intérieures s'annulent par orientation. Le lecteur
hydrodynamique interprète le membre de gauche comme une circulation et celui de droite
comme une vorticité accumulée. Le lecteur de jauge interprète le premier comme une
holonomie linéarisée et le second comme une courbure.

\begin{theorem}[Théorème du feuillet commun de bord]
\label{thm:hoja-comun-frontera}
La décomposition \(B_m=E_-B_m\oplus E_+B_m\) possède deux représentations
typées :
\[
\boxed{
\begin{array}{ccl}
E_-B_m&\longmapsto&
\text{circulation}\ /\ \text{holonomie},\\[2mm]
E_+B_m&\longmapsto&
\text{contrôle dissipatif}\ /\ \text{noyau réfléchi}.
\end{array}}
\]
L'égalité de Stokes et la coïncidence du spectre non nul appartiennent à
l'objet commun. La contraction hydrodynamique et l'écart spectral de jauge sont des conséquences
distinctes dans leurs codomaines respectifs.
\end{theorem}

\begin{proof}
Les projecteurs \(E_\pm\) sont complémentaires parce que
\(\theta_m^2=I\). L'identité de Stokes fournit l'entrelacement du
secteur orienté ; le \cref{thm:operador-comun-frontera-ns-ym} fournit le
spectre commun. Les
\cref{thm:contraccion-flujo-finito,thm:coercividad-gap-gauge-finito}
calculent les deux réalisations.
\end{proof}

Le théorème est une fermeture finie HMT : il n'identifie pas une EDP à une théorie
de jauge ni ne confond l'écart spectral \(3\) avec une masse physique. Il explique par construction
pourquoi la régularité intérieure et la séparation de bord sont deux fonctions de
la même mémoire orientée. Tout passage continu ultérieur doit conserver
l'opérateur, la norme et le spectre qui rendent cette identité exacte.
\endgroup
